
%%%%%%%%%%%%%%%%%%%%%%%%%%%%%%%%%%%%%%%%%%%%%%%%%%%%%%%%%%%%%
\subsection{问题三的模型的建立和求解}
\subsubsection{具体分析}

问题三要求在干扰源个数未知、位置未知、频道未知的条件下，为机器狗制定从原点出发、初始频道为 1 的自动搜索定位与清除策略，在确保全部清除的前提下最小化任务时间，属于不完全观测条件下的序贯决策问题。观测算子只返回角度且叠加 $1^\circ$ 误差，目标个数与位置先验完全未知，决策变量是"下一步去哪里、测哪个频道"。

问题的结构可分解为三个层次。第一层是探测的可达性，任何干扰源只有落在某个观测点的有效接收半径内才能被察觉，而有效接收半径在 $1000\sim1500$ 米之间且具体数值未知，因此必须设计对半径取最坏情形仍成立的观测点集。第二层是定位的欠定性，单次示向度只给出方向，必须至少两次不同位置的观测才能解算位置，而两次观测的几何关系直接决定定位精度。第三层是清除的距离约束，激光清除要求机器狗与干扰源的距离不超过 $20$ 米，因此定位精度必须优于该阈值。三个层次通过虚拟时间串联：探测次数消耗固定检测时间，定位机动消耗移动时间，逼近与清除消耗另一部分时间，策略设计本质上是在这三项之间做最优分配。

关于时间口径需要先行澄清。题目中"定位清除总时间包含机器狗移动时间、频道切换时间、检测时间、光学精确定位时间及清除时间"，这些量均由模拟器按虚拟世界时钟累计；而 $20$ 分钟的程序运行时间限制针对的是机器狗程序在现实世界中的运行时长。因此在虚拟时间上限 $360000$ 秒（$100$ 小时）之内，真正需要最小化的量是虚拟时间之和，现实时间限制仅要求程序本身足够高效。

\subsubsection{模型建立}

\paragraph{可检测性模型与保证可测判据} 由附录 2，全向干扰源 $G_c$ 在检测点 $P$ 处可被检测的充要条件为 $\lVert P-G_c\rVert\le r_c$，其中有效接收半径 $r_c\in[1000,1500]$ 且不通过接口返回。由于 $r_c$ 的取值未知，本文采用最坏情形口径：若观测点集 $\mathcal Q$ 满足
\begin{equation}
\forall X\in D,\quad \exists Q\in\mathcal Q:\ \lVert X-Q\rVert\le R_c^{\min}=1000,
\label{eq:cover-cond}
\end{equation}
则无论各干扰源的真实 $r_c$ 取何值，目标区域内的任何干扰源都至少能被 $\mathcal Q$ 中某一点测到。式 \eqref{eq:cover-cond} 把"保证可测"问题化归为纯粹的圆覆盖问题：用半径 $1000$ 米的圆盘覆盖半径 $1800$ 米的圆域 $D$。

\paragraph{观测点数下界与保证覆盖网} 覆盖问题 \eqref{eq:cover-cond} 的规模下界由 Kershner 的经典结论给出：$6$ 个单位圆盘所能覆盖的最大圆域半径恰为 $\sqrt3\approx1.73205$。将单位圆盘按半径 $1000$ 米缩放，则 $6$ 个半径 $1000$ 米的圆盘最多覆盖半径 $1732.05$ 米的圆域，而目标圆域半径为 $1800$ 米，故 $N_q\ge7$。取"$1$ 个中心点加 $6$ 个六边形环点、环半径 $d$"的对称构型，其覆盖半径为
\begin{equation}
\rho(d)=\max_{X\in D}\ \min_{Q\in\mathcal Q}\lVert X-Q\rVert .
\label{eq:rho-d}
\end{equation}
该式对 $d$ 先减后增：环半径过小时环点与中心点彼此靠得过近而圆周远端的覆盖被放弃，环半径过大时中心与环之间出现空档。$\rho(d)$ 的极小点 $d^{*}$ 给出该构型族的最优覆盖半径，而方程 $\rho(d)=R_c^{\min}$ 的两个根 $d_1<d^{*}<d_2$ 界定出可行的环半径区间。由于扫描行程等于六边形环的周长 $6d$，在可行区间内应取最小根 $d_1$ 以最小化行程，同时以 $\rho=1000$ 米恰好满足保证可测条件、不留多余裕量。

\paragraph{交会定位与逼近阈值} 由问题二的结论，对已知距离 $r$ 的干扰源，最优布站为 $b=\sqrt2r$ 且定位误差为 $E=3\delta r$。清除操作要求 $E\le20$ 米，由此反解出所需的逼近距离阈值
\begin{equation}
r_{\mathrm{thr}}=\frac{E_{\max}}{3\delta}=\frac{20}{3\times0.0174533}=381.97\ \text{m}.
\label{eq:threshold}
\end{equation}
式 \eqref{eq:threshold} 是连接问题二与问题三的关键桥梁：它把抽象的"定位精度"翻译为可执行的"逼近到 $382$ 米以内"。由于 $382$ 米远小于可行域尺度，单次远距离交会不足以直接定位，必须采用"粗定距加迭代逼近"的两阶段方案：先用较短侧移获得量级正确的距离估计，再沿视线方向逼近，最终由观测结果直接触发清除。

\paragraph{虚拟时间模型} 定义从出发到任务结束的虚拟总时间
\begin{equation}
T=\underbrace{\frac{L}{v}}_{\text{移动}}+\underbrace{5N_m}_{\text{检测}}
+\underbrace{N_{sw}}_{\text{频道切换}}
+\underbrace{5N_{c1}+3N_{c0}}_{\text{清除}},
\qquad v=5\ \text{m/s},
\label{eq:time-model}
\end{equation}
其中 $L$ 为累计位移，$N_m$ 为检测次数，$N_{sw}$ 为频道切换次数，$N_{c1},N_{c0}$ 分别为成功与未发现的清除次数。

时间模型 \eqref{eq:time-model} 可分解为两个可独立优化的子问题。其一是频道切换次数的最小化。在单点依次测量 $20$ 个频道时，若进入该点时当前频道为 $1$，则该点消耗 $19$ 次切换；在遍历多个观测点时，若每一点的测量顺序相对于上一点的末频道作反向安排，则每点仍为 $19$ 次切换，于是 $N_{sw}=N_q(N_c-1)=7\times19=133$。其二是在已知目标集合后规划访问顺序。清除阶段的位移可建模为以扫描终点为起点、以各干扰源为必访点的开放式旅行商问题。由 Beardwood--Halton--Hammersley 定理，均匀分布于面积 $A$ 区域内的 $N$ 个点的最优巡回长度几乎必然满足
\begin{equation}
L_{\mathrm{TSP}}\ \sim\ \beta\sqrt{NA},\qquad A=\pi R^2,\quad \beta\approx0.7124 .
\label{eq:bhh}
\end{equation}
代入 $R=1800$ 得 $A\approx1.018\times10^7\ \text{m}^2$，$N=10,13,16$ 时 $L_{\mathrm{TSP}}$ 分别约为 $7187$、$8195$、$9091$ 米。

综合以上各部分，问题三策略的虚拟总时间可写作
\begin{equation}
T_{\mathrm{model}}=\frac{6d_1}{v}+5N_cN_q+N_q(N_c-1)
+\frac{L_{\mathrm{TSP}}+N\beta_l\bar r}{v}
+N\bigl(2\cdot5+2\cdot1+5\bigr),
\label{eq:time-total}
\end{equation}
其中第三项为清除阶段的移动时间、$N\beta_l\bar r$ 为每个目标定位所需的侧移机动总量，第四项为每个目标在定位阶段消耗的两次补充检测、两次频道切换与一次成功清除。相应的理想下界（忽略全部信息获取开销，仅保留位移与最小动作）为
\begin{equation}
T_{\mathrm{lb}}=\frac{L_{\mathrm{TSP}}}{v}+(2\cdot5+1+5)N .
\label{eq:time-lb}
\end{equation}
式 \eqref{eq:time-lb} 的作用是提供一把标尺：任何实际策略的时间都不可能低于它，因此策略优劣可以用与下界的相对差距衡量，而不再依赖经验性评价。

\paragraph{分层搜索策略} 综合上述模型，本文给出四阶段闭环策略。第一阶段为全域粗探，机器狗自原点出发依次访问保证覆盖网点 $\mathcal Q=\{Q_0,\dots,Q_6\}$，其中 $Q_0$ 即出发点，在每个点依次对 $20$ 个频道执行检测并记录结果；由于式 \eqref{eq:cover-cond} 保证任何位于 $D$ 内的干扰源必被其中某点覆盖，此阶段结束时全部全向干扰源都至少会被登记一次示向度。第二阶段为粗定距，对每个已登记的频道，机器狗以当前位置为起点沿垂直于该示向度的方向侧移 $b=\beta\hat r$ 后复测，两条方位线交会给出粗略位置估计；侧移的目的不是直接达到清除精度，而是把径向不确定度从千米量级压缩到百米量级，为逼近提供方向。第三阶段为迭代逼近，机器狗沿当前视线方向朝估计位置直线行进，行进中每隔一段距离复测一次并更新位置估计；由于侧向角度误差与距离成正比，随着距离缩短单次方位观测的侧向定位误差按 $\lVert P-G\rVert\tan\delta$ 线性下降，逼近过程本身即可持续提高精度。第四阶段为就地清除，当检测结果返回 near（距离不超过 $5$ 米且位于覆盖角内）时直接执行清除，或当估计距离落入式 \eqref{eq:threshold} 给出的阈值以内时按交会结果定位后执行清除；清除成功后该频道被标记为已解决，在后续所有观测点不再复测，从而压缩后续检测次数。

\subsubsection{模型求解}

对式 \eqref{eq:rho-d} 做一维极小与求根，得到最优环半径 $d^{*}=1558.85$ 米对应最小覆盖半径 $\rho^{*}=900.00$ 米；方程 $\rho(d)=1000$ 的两个根为 $d_1=1122.96$ 米与 $d_2=1733.50$ 米。取 $d_1$ 得到最小行程的保证覆盖网，七点坐标如表 \ref{tab:q3-net} 所示，六边形环的周长为 $6d_1=6737.7$ 米。

\begin{table}[H]
	\centering
	\caption{七点保证覆盖网坐标（原点为出发点）}
	\label{tab:q3-net}
	\begin{tabular}{cccccccc}
		\toprule
		点编号 & $Q_0$ & $Q_1$ & $Q_2$ & $Q_3$ & $Q_4$ & $Q_5$ & $Q_6$ \\
		\midrule
		$x$ (m) & $0$ & $1122.96$ & $561.48$ & $-561.48$ & $-1122.96$ & $-561.48$ & $561.48$ \\
		$y$ (m) & $0$ & $0$ & $972.51$ & $972.51$ & $0$ & $-972.51$ & $-972.51$ \\
		\bottomrule
	\end{tabular}
\end{table}

对覆盖性做稠密数值验证：在目标圆域内均匀采样 $1.5\times10^6$ 个点并附加等分的边界采样，计算每点到最近观测点的距离并取最大值，结果为 $1000.000002$ 米，超出 $1000$ 米加 $10^{-6}$ 米数值容差的点仅 $2$ 个，覆盖率 $99.999867\%$，可认定为完全覆盖。同时本文还独立验证了下界 $N_q\ge7$：对纯 $6$ 个环点的构型做同样的覆盖半径优化，所得最小覆盖半径为 $1039.23$ 米，仍大于 $1000$ 米，与 Kershner 定理的断言一致，从两个方向确认了 $7$ 是必需的。

\begin{figure}[htbp]
	\centering
	\includegraphics[width=0.8\textwidth]{../求解/问题三/figures/七点保证覆盖网示意图.png}
	\caption{七点保证覆盖网与覆盖半径随环半径的变化}
	\label{fig:q3-net}
\end{figure}

如图 \ref{fig:q3-net}(a) 所示，七个半径 $1000$ 米的圆盘恰好覆盖整个目标圆域，相邻圆盘的交叠带较窄，说明该构型在覆盖意义上已接近最优；粗探巡回路径为六边形环加一次往返中心。图 \ref{fig:q3-net}(b) 给出覆盖半径随环半径的变化曲线，曲线在 $d^{*}=1558.85$ 米处触底于 $900$ 米，与横线 $1000$ 米相交于 $d_1$ 与 $d_2$ 两点，两点之间的阴影区即覆盖可行的环半径区间，本文取最左端的 $d_1$ 以获得最小行程。

由式 \eqref{eq:threshold}，进入清除精度所需的最大逼近距离为 $381.97$ 米，即当机器狗逼近到距源约 $382$ 米以内时，一次交会所得定位精度即可优于 $20$ 米的清除半径。

将覆盖网参数代入时间模型 \eqref{eq:time-total}，得到不同源个数下的虚拟时间分解，如表 \ref{tab:q3-time} 所示。

\begin{table}[H]
	\centering
	\caption{问题三虚拟时间模型与理论下界（侧移系数 $\beta_l=0.25$）}
	\label{tab:q3-time}
	\begin{tabular}{cccccc}
		\toprule
		干扰源个数 $N$ & 粗探时间 (s) & 清除移动 (s) & 清除动作 (s) & 总时间 (s) & 平均定位清除时间 (s) \\
		\midrule
		$10$ & $2180.5$ & $1796.8$ & $170.0$ & $4147.4$ & $414.7$ \\
		$11$ & $2180.5$ & $1884.6$ & $187.0$ & $4252.1$ & $386.6$ \\
		$12$ & $2180.5$ & $1968.3$ & $204.0$ & $4352.9$ & $362.7$ \\
		$13$ & $2180.5$ & $2048.7$ & $221.0$ & $4450.3$ & $342.3$ \\
		$14$ & $2180.5$ & $2126.1$ & $238.0$ & $4544.6$ & $324.6$ \\
		$15$ & $2180.5$ & $2200.7$ & $255.0$ & $4636.2$ & $309.1$ \\
		$16$ & $2180.5$ & $2272.9$ & $272.0$ & $4725.4$ & $295.3$ \\
		\midrule
		\multicolumn{6}{l}{理想下界 $T_{\mathrm{lb}}=L_{\mathrm{TSP}}/v+16N$：$N=13$ 时为 $1847.0$ s} \\
		\bottomrule
	\end{tabular}
\end{table}

表 \ref{tab:q3-time} 显示，粗探阶段固定耗时 $2180.5$ 秒，由 $6737.7$ 米巡回位移（$1347.5$ 秒）、$140$ 次检测（$700$ 秒）与 $133$ 次频道切换（$133$ 秒）构成，约占总时间的一半，是时间开销的主体；清除阶段的移动时间随 $N$ 线性增长，$N=13$ 时为 $2048.7$ 秒；清除阶段的动作时间仅占 $5\%$ 左右。总时间在 $N=10\sim16$ 区间内从 $4147$ 秒增至 $4725$ 秒，而平均定位清除时间从 $414.7$ 秒降至 $295.3$ 秒，后者随 $N$ 增大而减小，说明粗探成本被更多的目标分摊，干扰源越多反而"单位清除成本"越低。

\begin{figure}[htbp]
	\centering
	\includegraphics[width=0.8\textwidth]{../求解/问题三/figures/分层搜索几何结构图.png}
	\caption{分层搜索策略的几何结构（粗探、定距、逼近、清除）}
	\label{fig:q3-layer}
\end{figure}

如图 \ref{fig:q3-layer}(a) 所示，粗探阶段每个观测点以 $1000$ 米为半径覆盖一片区域，落入某片区域的干扰源会向该观测点提供一个示向度。图 \ref{fig:q3-layer}(b) 展示单个目标的定距与逼近过程：自粗探点沿垂直视线方向侧移到复测点，两条方位线交会得到带不确定度的位置估计，再沿视线逼近。图 \ref{fig:q3-layer}(c) 放大逼近末端，机器狗沿视线穿过干扰源附近，由于角误差随距离线性缩小，其横向偏差在接近干扰源时已小于 $5$ 米的近距离阈值，因而触发直接清除。

\begin{figure}[htbp]
	\centering
	\includegraphics[width=0.8\textwidth]{../求解/问题三/figures/虚拟时间构成与灵敏度图.png}
	\caption{虚拟时间构成、与理想下界的对比及侧移系数灵敏度}
	\label{fig:q3-time}
\end{figure}

如图 \ref{fig:q3-time}(a) 所示，粗探部分在各 $N$ 下均为固定常量，而清除阶段的移动时间随 $N$ 线性增长。图 \ref{fig:q3-time}(b) 把模型总时间与理想下界对比，$N=13$ 时前者为 $4450$ 秒、后者为 $1847$ 秒，前者约为后者的 $2.4$ 倍，差距主要来自粗探阶段不可避免的探测开销，说明本文策略在"硬保证全覆盖"这一约束下已属紧凑。图 \ref{fig:q3-time}(c) 显示侧移系数 $\beta$ 对平均定位清除时间的影响近似线性，$\beta$ 从 $0$ 增至 $0.5$ 时平均时间由 $297$ 秒增至 $366$ 秒，$23\%$ 的增幅对应着定位失败的完全规避，该折中是合理的。

按题目要求，问题三三次正式测试的结果应以规定格式给出，如表 \ref{tab:q3-test} 所示。

\begin{table}[H]
	\centering
	\caption{问题 3 正式测试结果}
	\label{tab:q3-test}
	\begin{tabular}{cccc}
		\toprule
		测试案例编码 & 清除干扰源个数 & 平均定位清除时间 (s) & 程序运行时间 (s) \\
		\midrule
		测试 1 案例编码 & & & \\
		测试 2 案例编码 & & & \\
		测试 3 案例编码 & & & \\
		\bottomrule
	\end{tabular}
\end{table}
